% 第 3 题程序框图：按固定源内容流矢量层精确值复刻（1 pt → 0.03515 cm，与全库同制）。
% 原点取框图竖轴 x=218.0000 pt、底缘 y（顶向下）=534.2764 pt；y 向上。
% 源实测：竖轴严格 x=218.0000；8 框（开始=圆角 27.999×16.838、x=1,y=1,k=0 矩形 94.904×16.333、
% s=x−y,t=x+y 矩形 98.847×15.499、x=s,y=t 矩形 61.655×15.499、k=k+1 矩形 50.677×15.166、
% k⩾3 菱形 74.733×18.560、输出(x,y) 平行四边形（上下边 62.304、右斜 6.449）、结束=圆角 27.999×17.016）。
% 连线：7 支竖直（x=218）+1 支回折（自菱形右顶点右行 22.678、上行 106.529、左行 56.740 入 N3 上方竖线）；
% 全图无虚线、无填充（除 8 支实心箭头）；框内文字 STSong/CM 系 10.5 pt，本图用 \small 同制。
\begin{tikzpicture}[x=0.03515cm,y=0.03515cm,>=stealth,every node/.style={font=\small,inner sep=0pt}]
% ===== 8 框（外轮廓照源实测值；圆角半径 3.985 单位=0.1401cm）=====
\draw[line width=0.398bp,rounded corners=0.1401cm] (-13.9993,217.9654) rectangle (13.9997,234.8034);   % N1 开始
\draw[line width=0.398bp] (-47.4520,187.0604) rectangle (47.4520,203.3934);                             % N2 x=1,y=1,k=0
\draw[line width=0.398bp] (-49.4240,156.9894) rectangle (49.4230,172.4884);                             % N3 s=x−y,t=x+y
\draw[line width=0.398bp] (-30.8270,126.9184) rectangle (30.8280,142.4174);                             % N4 x=s,y=t
\draw[line width=0.398bp] (-25.3380,97.1804) rectangle (25.3390,112.3464);                              % N5 k=k+1
\draw[line width=0.398bp] (37.3666,73.2460)--(0,82.5260)--(-37.3664,73.2460)--(0,63.9660)--cycle;       % N6 菱形 k⩾3
\draw[line width=0.398bp] (-27.9312,49.3093)--(34.3728,49.3093)--(27.9238,31.5893)--(-34.3802,31.5893)--cycle; % N7 输出 (x,y)
\draw[line width=0.398bp,rounded corners=0.1401cm] (-13.9993,0) rectangle (13.9997,17.0160);            % N8 结束
% ===== 连线（7 竖 + 1 回折；箭头顶端照源落在下一框边线上）=====
\draw[->,line width=0.398bp] (0,217.7662)--(0,203.3934);                        % L1 开始→x=1,y=1,k=0
\draw[->,line width=0.398bp] (0,187.0604)--(0,172.4884);                        % L2 →s=x−y,t=x+y（回折靶线）
\draw[->,line width=0.398bp] (0,156.9894)--(0,142.4174);                        % L3 →x=s,y=t
\draw[->,line width=0.398bp] (0,126.9184)--(0,112.3464);                        % L4 →k=k+1
\draw[->,line width=0.398bp] (0,97.1804)--(0,82.5260);                          % L5 →菱形尖顶
\draw[->,line width=0.398bp] (0,63.9660)--(0,49.3093);                          % L6 →输出
\draw[->,line width=0.398bp] (0,31.5893)--(0,17.0160);                          % L7 →结束
\draw[->,line width=0.398bp] (37.3666,73.2460)--(60.3263,73.2460)--(60.3263,179.7744)--(0,179.7744); % L8 否支回折
% ===== 框内文字与分支标签（照源字形中心定位）=====
\node at (0,226.9660) {开始};
\node at (0,195.2264) {\(x=1,\ y=1,\ k=0\)};
\node at (0,165.1564) {\(s=x-y,\ t=x+y\)};
\node at (0,135.0864) {\(x=s,\ y=t\)};
\node at (0,104.1764) {\(k=k+1\)};
\node at (0,72.9414) {\(k\geqslant 3\)};
\node at (0,41.2114) {输出 \((x, y)\)};
\node at (0,9.1164) {结束};
\node at (8.9500,57.2864) {是};
\node at (37.6500,82.2164) {否};
\end{tikzpicture}
